\SourceLeaf{154}
S. [128 Z. 29]; der Zinsfuss aber, obwohl von seiner frühern Höhe sehr herabgegangen, noch immer hoch genug, dass die Kasse jener höhern Pensionenzahl noch gewachsen bleibt. Überhaupt ist diese hohe Zahl, selbst wenn sie noch um eine oder ein paar Pensionen mehr gestiegen wäre, an sich, und in so weit man darin nur das Vorkommen eines ungewöhnlich hohen Schwankens über den Mittelwerth 15 oder 17 zu erkennen hat, lange nicht von einer so schweren Bedeutung, wie \textit{der} Umstand, zu welchem ich jetzt übergehe
\begin{quote}
dass der Mittelwerth der Pensionenzahl, möge man 15 oder 17 wie den plausibelsten betrachten, nur so lange gültig ist, als die Genossenschaft keine grössere Ausdehnung erhält, als sie 1794 hatte, und dass diese Gültigkeit jetzt, wo die Ausdehnung, mit dem Maassstabe der Zahl der verheiratheten Mitglieder gemessen, um mehr als 60 Procent grösser ist, als zur Zeit jener Epoche, ganz aufgehört hat.
\end{quote}
\noindent\hspace*{\parindent}Für die Zwischenzeit zwischen 1794 und 1845 lässt sich dieser Maassstab nicht anwenden, weil die Kenntniss der Zahl der verheiratheten Mitglieder fehlt. So viel sich aber aus der Bewegung der Gesammtzahl aller Mitglieder schliessen lässt, ist die Ausdehnung der Genossenschaft im Ganzen und abgesehen von einigem hin und her Schwanken bis etwa 1831 nicht grösser, sondern eher etwas geringer gewesen als 1794, und \textit{sehr bedeutend} ist die Vergrösserung erst seit wenigen Jahren geworden. Mit der hohen Witwenzahl in der letzten Zeit steht daher die erweiterte Ausdehnung der Genossenschaft durchaus nicht in ursächlichem Zusammenhange, wie bereits oben [S. 128] bemerkt ist.

Der vierte hier noch zu betrachtende Punkt und gleichsam der Schlussstein der ersten Abtheilung, ist, einen Überschlag der künftigen Bewegung der Witwenzahl im Allgemeinen zu machen, so weit nemlich ein solcher auf dem bisher eingeschlagenen Wege erreicht werden kann. Die Anzahl aller Professoren ist jetzt 57; davon nehmen an der Witwenkasse Theil 51, und unter diesen sind 42 verehelicht. Macht man, zuerst, den Überschlag nach der Hypothese, die den Verhältnissen der hiesigen Professorenwitwenkasse am meisten angemessen scheint, dass auf 7 stehende Ehen 4 Witwen zu rechnen sind, so gibt die Rechnung 24 Witwen, und die Bedeutung davon ist, dass nachdem die Gesellschaft von jenem Umfange den Beharrungszustand erreicht hat, die durchschnittliche Witwenzahl 24 sein wird. Wegen der Waisenpensionen müsste nach \textsc{Kritter}’s Gutachten noch ein Sechstel zugesetzt werden; ich will jedoch nur ein Achtel in Rechnung bringen, also die durchschnittliche Pensionenzahl im Beharrungszustande \(=27\) setzen. Allerdings soll man den Beharrungszustand erst nach 45—50 Jahren erwarten: man würde sich aber sehr täuschen, wenn man diese weite Entfernung für einen starken Beruhigungsgrund hielte. Denn schon lange vorher ist man dem Grenzzustande so nahe gekommen, dass der Unterschied nicht viel mehr bedeutet. Man vergleiche die oben S. [141] für allmähliges Steigen der Witwenzahl mitgetheilten Ansätze, die ohne Anspruch auf strenge Genauigkeit zu machen, doch einigermassen eine Idee von dem Hergange geben können. Wenn man dann dabei überlegt, dass der frühere Umfang der Gesellschaft schon 17 Pensionen als Normalzahl (bei dem Verhältniss \(7:4\)) für die Pensionen gegeben hatte, und es also hier sich nur um die allmählige Entwicklung der auf 10 angeschlagenen Vergrösserung handelt, so wird man mich leicht verstehen, wenn ich behaupten muss, dass schon nach etwa 25 Jahren man auf 25 Pensionen als \textit{Mittelzahl} gefasst sein müsse. Hierzu kommt nun noch die Erweiterung des Spielraumes wegen der Schwankungen, denen scharfe Grenzen zu setzen unmöglich ist. Soviel ist aber gewiss, dass ein Schwanken von 7 Pensionen, auf und ab vom Mittelwerthe, wie etwas gar nicht ausserordentliches in den Überschlag mit aufgenommen werden muss, da ein verhältnissmässig wenigstens eben so grosses Schwanken nach frühern Präcedentien factisch ist. Das Resultat dieser Erwägungen ist also, dass man erwarten muss, um das Jahr 1870 die Pensionenzahl zwischen 18 und 32 zu finden, ohne dass man im Stande ist, im voraus zu bestimmen, \textit{wo} innerhalb dieses weiten Spielraumes; dass das Erreichen des einen oder des andern Extrems nicht wie etwas sehr ausserordentliches betrachtet werden darf; endlich, dass späterhin diese Zahlen noch ein
